Nous avons utilisé des variables comme compteurs de boucles et dans les expressions. 
L'instruction d'affectation
permet de les utiliser également pour mémoriser des résultats intermédiaires dans les
calculs.

\section{Variables et affectation}
Une variable\index{Variable} est utilisée pour conserver une valeur d'un certain type dans la mémoire
de l'ordinateur. Une déclaration de variable\index{Déclaration de variable} indique le nom de la variable (qui est un identificateur)
et son type~:

\vspace{\neglskip}\mbox{\blc  \blk\idt  \mcl var\idt }\par 
\vspace{\neglskip}\mbox{\blc  \blk\blk\idt  B\ $:$\ boolean$;$}\par 
\vspace{\neglskip}\mbox{\blc  \blk\blk\idt  I\ $:$\ integer$;$}\par 
\vspace{\neglskip}\mbox{\blc  \blk\blk\idt  R\ $:$\ real$;$}\par 
\vspace{\neglskip}\mbox{\blc  \blk\blk\idt  S\ $:$\ string$;$}\par

\noindent C'est le compilateur qui décide quelle partie de la mémoire et quel codage sera utilisé pour
représenter la valeur de la variable. Comme la mémoire contient toujours quelque chose (puisque
chaque bit qui la constitue vaut 0 ou 1), la
variable a toujours une valeur. On peut utiliser cette valeur dans une expression (comme
$((2\:\ast$ {\idt I} $) + 1)$) en la notant
simplement par le nom de la variable ({\idt I} dans l'exemple). Au début de l'exécution du programme, la valeur
d'une variable est {\em indéterminée} en ce sens qu'elle correspond à l'état initial de la 
mémoire qui peut être quelconque.  L'instruction d'{\em affectation\/} \index{Affectation}permet de changer la 
valeur d'une variable. Par exemple l'instruction~:

\vspace{\neglskip}\mbox{\blc  \blk\blk\idt  B\ $:=$\ true$;$}\par 

\noindent affecte la valeur ``vrai'' à la variable booléenne {\idt B}. De manière générale, une
instruction d'affectation a la forme suivante~:

\index{Affectation}
\begin{center}
\begin{tabular}{@{}|p{12.5cm}|@{}} \hline
{\em Instruction d'affectation~:}\\ \hline \hline
\begin{minipage}[t]{12.5cm}
{\medskip

\vspace{\neglskip}\mbox{\blc  \blk\blk\idt  NomDeVariable\ $:=$\ \rm Expression%
\footnote{\ Le type de l'expression est {\em compatible} avec celui de la variable.}$;$}\par 

\medskip}
\end{minipage}\\ \hline 
\end{tabular}
\end{center}

\noindent Le symbole {\idt  $:=$} se lit ``devient égal à''. L'exécution de cette instruction consiste à
évaluer l'``Expression'' pour calculer sa valeur qui devient celle de la variable nommée
{\idt  NomDeVariable}. Le type de l'expression est compatible avec celui de la variable dans les cas 
suivants~:

\index{Compatibilité de type}
\begin{center}
\begin{minipage}[t]{8.2cm}
\begin{tabular}{@{}|l|l|@{}} \hline
\multicolumn{2}{@{}|l|@{}}{\em Compatibilité de type pour l'affectation~:} \\ \hline \hline
Type de la variable & Type de l'expression\\ \hline
boolean & boolean \\
integer & integer \\
integer & longint\/$^{b}$ \\
longint & integer \\
longint & longint \\
real & integer \\
real & longint \\
real & real \\
string & string \\ 
\hline
%\multicolumn{2}{@{}|l|@{}}{\noindent \rule{33mm}{0.1mm}} \\
\multicolumn{2}{@{}|l|@{}}{\rule{0mm}{4.5mm}$^{b}$\ \footnotesize C'est \hfill une \hfill erreur \hfill à \hfill l'exécution \hfill si \hfill la \hfill valeur \hfill de } \\
\multicolumn{2}{@{}|l|@{}}{\footnotesize l'expression \hfill n'est \hfill pas \hfill comprise \hfill entre \hfill {\idt $-$maxint} \hfill et } \\
\multicolumn{2}{@{}|l|@{}}{\footnotesize {\idt maxint}.} \\ \hline 
\end{tabular}
\end{minipage}
\end{center}

%footnote{\ C'est une erreur à l'exécution si la valeur de l'expression 
%n'est pas comprise entre {\idt $-$maxint} et {\idt maxint}.}

\noindent  Quand les types de l'expression et de la variable sont 
compatibles mais ne sont pas les mêmes, le compilateur 
effectue des {\em conversions}\index{Conversion}
c'est-à-dire les changements de codage appropriés.

\section{Sauvegarde de l'état du robot}

Une première utilisation des variables est la sauvegarde de l'état du \robot. Imaginons par 
exemple que nous voulions concevoir une procédure pour faire un dessin qui nécessite divers 
changements de la taille de la grille. Pour que l'appel de
cette procédure dans un programme ne modifie pas la taille de la grille utilisée dans le 
programme, il faut sauvegarder cette
taille au début de la procédure puis la restituer à la fin, comme dans l'exemple suivant~:

{\raggedright 
 %\vspace{\baselineskip}
 \vbox to 0pt{\includevspacePostScript[13mm]{Canards}{0.95}{0.95}}%
 \input Canards.tex
 }
 
\section{Compter}
Supposons que nous voulions
compter le nombre d'itérations dans une boucle ``while''. On utilise une variable déclarée de 
type entier~: 

\mbox{\blc  \blk\idt  \mcl var\idt\ I\ $:$\ integer$;$}\par 

\noindent Avant d'exécuter la boucle ``while'', la variable doit être initialisée à 0~:

\mbox{\blc  \blk\blk\idt  I\ $:=$\ 0$;$}\par 

\noindent Après chaque tour de boucle, la valeur de la variable doit être augmentée de 1~:

\mbox{\blc  \blk\blk\idt  I\ $:=$\ $($I\ $+$\ 1$);$}\par 

\noindent Supposons que la valeur de la variable {\idt I} soit égale à 5 {\em avant} l'affectation 
\mbox{\idt  I\ $:=$\ $($I\ $+$\ 1$);$}. La valeur de l'expression {\idt  $($I\ $+$\ 1$)$}
est $(5 + 1) = 6$ de sorte que la valeur de la variable {\idt I} {\em après} l'affectation sera égale
à 6.

Par exemple, ceci peut nous servir pour déterminer la taille de la fenêtre dans laquelle le \robot\ 
se déplace sur l'écran de l'ordinateur~:
\input CalculerTailleFenetreRobot.tex

\section{Entrée d'une valeur d'une variable}
On peut changer la valeur d'une variable en entrant sa nouvelle valeur au clavier de 
l'ordinateur. Pour
ce faire la commande du \robot\ {\idt modeTexte$;$} fait apparaître la fenêtre prévue pour les
entrées-sorties en \pascal\ (cette fenêtre est cachée par la fenêtre 
de déplacement du \robot). Ensuite 
l'instruction {\idt readln$($X$);$} permet d'entrer la valeur de la variable {\idt X} au 
clavier, cette valeur s'inscrivant en même temps dans la fenêtre d'entrées-sorties de \pascal.
La procédure {\idt readln$($X$);$} est {\em polymorphe}\index{Polymorphe} ce qui veut dire que le type de la va\-ria\-ble {\idt X} peut
être {\idt boolean}, {\idt integer}, {\idt longint}, {\idt real} ou {\idt string} (alors que
la plupart des autres procédures en \pascal\ sont {\em monomorphes}\index{Monomorphe} en ce sens que leurs paramètres
ont toujours le même type\,\footnote{\ Le type du paramètre formel et le type du
paramètre effectif doivent être compatibles pour l'affectation.}). La valeur que l'on tape au clavier s'écrit comme les constantes du type
de {\idt X} sauf que pour les chaînes de caractères, il ne faut pas utiliser 
 d'apostrophes. On peut 
corriger les fautes de frappe avec la touche  %touche d'effacement 
\framebox[11.mm][l]{$\leftarrow$~~~}\hspace{0.3ex} 
d'effacement. Il faut terminer l'entrée de la valeur par la frappe de la touche %retour a la ligne
\framebox[11.mm][l]{$\hookleftarrow$~~~}\hspace{0.3ex}
de retour à la ligne.

L'instruction {\idt readln$;$} (sans paramètre) permet d'attendre que la touche %retour a la ligne
\framebox[11.mm][l]{$\hookleftarrow$~~~}\hspace{0.3ex}
de retour à la ligne soit enfoncée. Elle est donc utile pour éviter que la fenêtre prévue pour les
entrées-sorties en \pascal\ ne disparaisse à la fin de l'exécution du programme ou lorsque
l'exécution d'une commande quelconque du \robot\ fait réapparaître 
la fenêtre de dessin. 

\section{Affichage de la valeur d'une expression}
Pour afficher la valeur d'une expression $E$ dans la fenêtre d'entrées-sorties de \pascal, on
peut utiliser l'instruction polymorphe {\idt write$(E);$}. Les écritures se font sur la même ligne, les unes 
à la suite des autres, en passant à la ligne suivante en fin de ligne. L'instruction
{\idt writeln$;$} permet de forcer le passage à la ligne. Par exemple 
le programme ci-dessous lit
un entier $n$ 
puis écrit sa factorielle $n! = (1 \ast 2 \ast \ldots \ast (n - 1) \ast n)$~:
\input Factorielle.tex

\vfill

\noindent \rule{137mm}{0.1mm} %---------------------------------------------

\noindent\begin{exercices}{\ \\

\tiret Écrire un programme qui dessine un cadre le plus près possible du bord de la fenêtre de déplacement
du \robot.

\tiret Écrire un programme qui esquisse les dessins à  compléter de la figure \ref{fig:completer-les-dessins} 
(page \pageref{fig:completer-les-dessins}).

\tiret La numération babylonienne \cite{guitel}, qui date d'environ 2000 ans av. J.-C., est en
base 60. C'est donc une numération sexagésimale tout comme celle que nous utilisons encore
pour noter l'heure. Les unités, de 1 à 59, sont notées de manière additive en accumulant des ``clous'' (pour
1) et des ``chevrons'' (pour 10) regroupés et superposés comme indiqué sur la figure
ci-dessous.\\[3.9cm]
\noindent\includePostScript{Babyloniens}{0.95}{0.95}%[3.9cm]

\noindent Au delà de 60, la notation est positionnelle, la valeur d'un chiffre occupant
la $p^{\grave{e}me}$ position (comptée de droite à gauche à partir de zéro) étant
multipliée par $60^{p}$. L'absence d'unités sexagésimales en position médiale (que
nous no\-te\-rions par 0) fut d'abord marquée par un espace puis, vers le $3^{\grave{e}me}$
siècle av. J.-C., par ``deux clous inclinés''. Les Babyloniens ont donc mis 15 siècles pour
inventer une notation pour le
zéro sans toutefois en comprendre la valeur calculatoire$\:$! Par exemple 
le nombre $3657 = 1 \times 60^{2} + 0 \times 60^{1} + 57 \times 60^{0}$
est représenté par la juxtaposition des chiffres $1$, $0$ et 
$57 = 50 + 7$~:

\pagebreak
%\vspace{-6mm}
\mbox{}\\[1.9cm]
\noindent\includePostScript{Babyloniens3657}{0.95}{0.95}%[2.4cm]

Écrire un programme \pascal\ de conversion
d'un nombre po\-si\-tif quel\-con\-que dans la nu\-mé\-ra\-tion babylonienne. 

\tiret La numération égyptienne ($3^{\grave{e}me}$ 
millénaire av. J.-C.) est décimale et additive \cite{guitel}. Elle utilise des hiéroglyphes pour 
noter les
puissances de 10 (le trait vertical pour 1, l'anse pour 10, la spirale pour 100, le lotus pour 1000,
l'index pour 10~000,
le têtard pour 100 000 et le génie pour 1 000 000)~:\\[3.0cm]

\noindent\includePostScript{ChiffresEgyptiens}{0.95}{0.95}%[3.0cm]
 
\noindent La valeur d'une séquence de signes est la somme des valeurs
de ces signes~:\\[2.0cm]

\noindent\includePostScript{1234567Egyptien}{0.95}{0.95}%[3.0cm]
 
\noindent On écrit de gauche à droite ou de droite à gauche, la
lecture se faisant dans le sens où sont orientés les hiéroglyphes. 
Les signes identiques se regroupent sur une, deux ou trois lignes de
la manière suivante~:\\[0.8cm]
 
%\vspace{-\baselineskip}
\noindent\includePostScript{GroupesEgyptiens}{0.95}{0.95}%[1.8cm]
 
\noindent  Entre -2000 et -1600 ans av. J.-C., une représentation hybride (additive et partiellement
positionnelle) est apparue. Un nombre inférieur à 1000 situé avant, au dessus ou au dessous
du signe du génie du million, du têtard, de l'index ou du lotus représente le produit de ce nombre et de la
valeur de ce hiéroglyphe. Par exemple, le nombre 37 280 549 peut se décomposer en
$((372 \ast 100 000) + 80 549)$ et s'écrit~:\\[2.3cm]
 
\noindent\includePostScript{37280549Egyptien}{0.95}{0.95}
 
\noindent Écrire un programme \pascal\ de conversion d'un nombre dans la nu\-mé\-ra\-tion
é\-gyp\-tien\-ne.

\tiret La numération acrophonique grecque \cite{guitel} fut utilisée à partir 
du V$\grave{e}me$
siècle av. J.-C. à Athènes. C'est une numération additive décimale utilisant la base
auxiliaire 5 pour diminuer les répétitions. Les chiffres supérieurs
à un sont notés par des 
lettres (ou combinaisons de lettres) 
correspondant chacune à l'initiale (ou à l'abréviation) du nom du
nombre correspondant~:\\[1.8cm]

\noindent\includePostScript{ChiffresGrecques}{0.95}{0.95}

Écrire un programme \pascal\ de conversion d'un nombre in\-fé\-rieur à 100~000 dans la nu\-mé\-ra\-tion
grecque. 
Pour représenter les grands nombres, utiliser la numération acrophonique 
pour les monnaies~:

\noindent Oboles~:\\[0mm]

%\vspace{-0.5cm}
\noindent\includePostScript{ObolesGrecques}{0.95}{0.95}

\noindent Drachmes (1 drachme = 6 oboles)~:\\[0.8cm]

\noindent\includePostScript{DrachmesGrecques}{0.95}{0.95}

\noindent Talents (1 talent = 6~000 drachmes)~:\\[0.8cm]

\noindent\includePostScript{TalentsGrecques}{0.95}{0.95}

\clearpage
\noindent Par exemple le nombre 239~686~966 se représente par~:\\[1.0cm]

\noindent\includePostScript{239686966Grecque}{0.95}{0.95}

\tiret La numération romaine \cite{guitel} est additive et soustractive (tout signe placé à gauche
d'un signe de valeur supérieure s'en soustrait). Ses chiffres
sont bien connus~:\\[1.2cm]

\noindent\includePostScript{ChiffresRomains}{0.95}{0.95}

Pour les grands nombres, la barre horizontale multiplie par 1~000 la valeur du nombre qu'elle
surligne et le rectangle sans base multiplie par 100~000 la valeur du nombre qu'il enferme,
comme dans l'exemple 135~792~468 ci-dessous~:\\[1.4cm]

\noindent\includePostScript{135792468Romain}{0.95}{0.95}

Écrire un programme \pascal\ de conversion d'un nombre dans la nu\-mé\-ra\-tion
ro\-mai\-ne.

\tiret Écrire un programme de dessin de la courbe de Sierpinski à un ordre
quelconque de la figure \ref{fig:CourbSierp} (page \pageref{fig:CourbSierp}).

\tiret Un carré magique d'ordre $n$ est un tableau carré de $n^{2}$ nombres
tel que la somme des nombres de chaque ligne, de chaque colonne et de chaque
dia\-go\-na\-le est la même. Des exemples sont donnés à la figure \ref{fig:Cmagic}
(page \pageref{fig:Cmagic}). Notons $C_{i, j}$ le nombre figurant à la
$i^{\grave{e}me}$ ligne et la $j^{\grave{e}me}$ colonne du carré comptées à
partir de 1 et numérotées de bas en haut pour les lignes et de gauche à
droite pour les colonnes. Écrire un programme \pascal\ qui cons\-truit des
carrés magiques d'ordre $n$ impair en commençant par placer 1 dans
$C_{i, j} = C_{\frac{n + 1}{2}, n}$ puis affecte les nombres 2, 3, \ldots, $n^{2}$
aux éléments $C_{i, j}$ obtenus en augmentant $i$ et $j$ de 1 modulo $n$ à chaque fois
pendant $n - 1$ pas, et en diminuant $j$ de 1 tout en laissant $i$ inchangé à chaque 
$n^{\grave{e}me}$ pas.

\noindent $\Box$}\end{exercices}
\begin{figure}[htbp]
 \includevspacePostScript[9.2cm]{CourbeSierpinski0a5}{0.95}{0.95}
 \caption{Courbes de Sierpinski d'ordres 0 à 5}\label{fig:CourbSierp}
\end{figure}
\begin{figure}[htbp]
 \includevspacePostScript[6.0cm]{CarreMagique}{0.95}{0.95}
 \caption{Carrés magiques d'ordres 3, 5 et 7}\label{fig:Cmagic}
\end{figure}

\pagebreak

\begin{corrige}{
Nous donnons un programme, les autres se trouvant sur disquette~:
\input CarreMagique.tex
\noindent $\Box$}\end{corrige}

%%%%%%%%%%%%%%%%%%%%%%%%%%%%%%%%%%%%%%%%%%%%%%%%%%%%%%%%%%%%%%%%%%%%%%%%%%%
\noindent \framebox[\textwidth]{
\begin{minipage}{\myparboxwidth}
\setlength{\parindent}{2em}
\small \sf
\vspace{\abovedisplayskip}
\section*{Addenda}
Une \underline{variable}\index{Variable} sert à donner un {\em nom} à une mémoire qui conserve 
une {\em valeur} représentée
selon un code déterminé par le {\em type} de la variable. Pour un codage binaire, on
peut imaginer que la mémoire est une suite de lampes éteintes (pour représenter
le bit 0) ou allumées (pour le bit 1).
La valeur
initiale de la va\-ria\-ble est {\em indéterminée} puisqu'elle dépend de l'état
de la mémoire au moment du début de l'exécution du programme (certains compilateurs
l'initialisent à zéro). L'instruction d'\underline{affectation}\index{Affectation} permet de 
conserver une nouvelle valeur en mo\-di\-fiant physiquement
l'état de cette mémoire. La nouvelle valeur est donnée par une expression (qui peut
utiliser l'ancienne valeur). La notion informatique de variable est donc très différente
de la notion mathématique homonyme. En mathématiques une variable désigne une valeur 
inconnue fixée une fois pour toutes. Pour un mathématicien, une variable informatique
serait la suite des valeurs qu'elle prend à sa création puis après chaque
affectation.

\vspace{\baselineskip}

On peut déclarer des variables {\em locales} à une procédure (ou fonction),
en insérant la déclaration de la variable entre l'en-tête et le corps.
Dans ce cas la variable n'est utilisable que dans le corps de la procédure. 
À l'exécution, elle n'existe (en ce sens qu'un espace
mémoire lui est alloué) que pendant la durée de l'exécution du corps
de la procédure. Pour une procédure récursive,
il s'agit à chaque appel d'une nouvelle variable sans aucun
lien avec les instances correspondant aux appels récursifs antérieurs.

Dans le corps de la procédure on peut également utiliser des variables
{\em glo\-ba\-les} dont la déclaration figure dans le programme
avant celle de la procédure. Ces variables globales existent
pendant toute la durée de l'exécution du programme. Par conséquent en cas
d'appels récursifs d'une procédure, il s'agit toujours de la
même variable globale.

Enfin des variables locales et globales peuvent porter le même nom. Dans
ce cas c'est la variable locale qui est visible dans la procédure, la variable
globale continuant à exister mais étant inaccessible.

\vspace{\abovedisplayskip}
\end{minipage}
}
%%%%%%%%%%%%%%%%%%%%%%%%%%%%%%%%%%%%%%%%%%%%%%%%%%%%%%%%%%%%%%%%%%%%%%%%%%%% 

